\heading{理论物理基础（II）-模拟题1-期末}
\noteinfo{题目和答案由AI生成。}

\begin{question}
考虑半径为 $R$ 的圆环上的非相对论粒子，其角坐标记为 $\phi\in[0,2\pi)$，满足周期性边界条件 $\psi(\phi+2\pi)=\psi(\phi)$。哈密顿量为
\[
\hat H=-\frac{\hbar^2}{2mR^2}\frac{d^2}{d\phi^2}.
\]

\begin{enumerate}[label=(\arabic*),leftmargin=2.5em]
  \item 求 $\hat H$ 的全部本征值与归一化本征态；
  \item 给定初态
  \[
  \psi(\phi)=A(1+\cos\phi),
  \]
  求归一化常数 $A$，并将其展开到能量本征态基底；
  \item 在该初态下测量能量，求所有可能的测量结果与概率；
  \item 设系统按 $\hat H$ 演化，是否存在 $t>0$ 使得 $\psi(\phi,t)=\psi(\phi,0)$？若有，求出所有这样的时刻；是否存在 $t>0$ 使得 $\psi(\phi,t)=-\psi(\phi,0)$？
\end{enumerate}

\begin{solution}
在圆环上，定态方程解为平面波 $e^{in\phi}$。周期性条件要求 $n\in\mathbb Z$。因此
\ansmath{\psi_n(\phi)=\frac{1}{\sqrt{2\pi}}e^{in\phi},
\qquad
E_n=\frac{\hbar^2 n^2}{2mR^2},
\qquad n\in\mathbb Z.}

对初态，归一化条件给出
\[
1=A^2\int_0^{2\pi}(1+\cos\phi)^2d\phi=A^2(2\pi+\pi)=3\pi A^2,
\]
故
\[
A=\frac{1}{\sqrt{3\pi}}.
\]
又因
\[
1+\cos\phi=1+\frac{e^{i\phi}+e^{-i\phi}}{2},
\]
所以
\[
\psi(\phi)=\sqrt{\frac{2}{3}}\,\psi_0(\phi)+\frac{1}{\sqrt6}\,\psi_1(\phi)+\frac{1}{\sqrt6}\,\psi_{-1}(\phi).
\]
即
\ansmath{A=\frac{1}{\sqrt{3\pi}},
\qquad
\psi=\sqrt{\frac23}\,\psi_0+\frac{1}{\sqrt6}\,\psi_1+\frac{1}{\sqrt6}\,\psi_{-1}.}

因此测量能量时，可能结果只有
\[
E_0=0,
\qquad
E_{\pm1}=\frac{\hbar^2}{2mR^2}.
\]
注意 $n=\pm1$ 对应同一个能量。概率为
\ansmath{P(E=0)=\frac23,
\qquad
P\!\left(E=\frac{\hbar^2}{2mR^2}\right)=\frac13.}

随时间演化后
\[
\psi(\phi,t)=\sqrt{\frac23}\,\psi_0(\phi)+e^{-iE_1t/\hbar}\left[\frac{1}{\sqrt6}\,\psi_1(\phi)+\frac{1}{\sqrt6}\,\psi_{-1}(\phi)\right],
\qquad E_1=\frac{\hbar^2}{2mR^2}.
\]
要使其回到初态，只需满足
\[
e^{-iE_1 t/\hbar}=1,
\]
故
\ansmath{t=\ell\,\frac{4\pi mR^2}{\hbar},\qquad \ell=1,2,3,\dots.}

而若要满足 $\psi(\phi,t)=-\psi(\phi,0)$，则 $n=0$ 分量也必须变号，但它始终不随时间改变，因此这是不可能的。故
\anstext{不存在 $t>0$ 使得 $\psi(\phi,t)=-\psi(\phi,0)$。}
\end{solution}
\end{question}

\begin{question}
考虑一个自旋为 $1$ 的粒子。取 $\hat S_z$ 的本征态 $\{|1\rangle,|0\rangle,|-1\rangle\}$ 为基。定义哈密顿量为
\[
\hat H=D\hat S_y,
\]
其中 $D$ 为实常数。

\begin{enumerate}[label=(\arabic*),leftmargin=2.5em]
  \item 求算符 $\hat S_y$ 在 $\{|1\rangle,|0\rangle,|-1\rangle\}$ 基下的矩阵表示；
  \item 求 $\hat S_y$ 的本征值与一组归一化本征态；
  \item 若初态为 $|1\rangle$，求时刻 $t$ 的量子态 $|\psi(t)\rangle$；
  \item 在态 $|\psi(t)\rangle$ 下先测量 $\hat S_z$，再测量 $\hat S_y$，求所有可能测量结果的组合及其概率。
\end{enumerate}

\begin{solution}
自旋 $1$ 的角动量算符满足
\[
\hat S_y=\frac{1}{\sqrt2}
\begin{pmatrix}
0&-i\hbar&0\\
i\hbar&0&-i\hbar\\
0&i\hbar&0
\end{pmatrix}
\]
（基的顺序为 $|1\rangle,|0\rangle,|-1\rangle$）。因此
\ansmath{\hat S_y\ \longleftrightarrow\ \frac{\hbar}{\sqrt2}
\begin{pmatrix}
0&-i&0\\
i&0&-i\\
0&i&0
\end{pmatrix}.}

$\hat S_y$ 的本征值为 $+\hbar,0,-\hbar$。可取归一化本征态为
\[
|+\rangle_y=\frac{1}{2}
\begin{pmatrix}-1\\-i\sqrt2\\1\end{pmatrix},
\qquad
|0\rangle_y=\frac{1}{\sqrt2}
\begin{pmatrix}1\\0\\1\end{pmatrix},
\qquad
|-\rangle_y=\frac{1}{2}
\begin{pmatrix}-1\\i\sqrt2\\1\end{pmatrix}.
\]
也就是
\ansmath{\hat S_y|\pm\rangle_y=\pm\hbar|\pm\rangle_y,
\qquad
\hat S_y|0\rangle_y=0.}

将初态 $|1\rangle$ 按 $\hat S_y$ 本征态展开后再演化，或直接利用旋转矩阵，都可得
\[
|\psi(t)\rangle=e^{-iDt\hat S_y/\hbar}|1\rangle
=\cos^2\frac{Dt}{2}\,|1\rangle+\frac{\sin Dt}{\sqrt2}|0\rangle+\sin^2\frac{Dt}{2}\,|-1\rangle.
\]
因此
\ansmath{|\psi(t)\rangle=
\cos^2\frac{Dt}{2}\,|1\rangle+\frac{\sin Dt}{\sqrt2}|0\rangle+\sin^2\frac{Dt}{2}\,|-1\rangle.}

第一次测量 $\hat S_z$ 时，可能结果与概率分别为
\[
S_z=+\hbar,\\ P=\cos^4\frac{Dt}{2};
\qquad
S_z=0,\\ P=\frac12\sin^2 Dt;
\qquad
S_z=-\hbar,\\ P=\sin^4\frac{Dt}{2}.
\]
坍缩后分别为 $|1\rangle,|0\rangle,|-1\rangle$。再测量 $\hat S_y$ 时：
\[
|1\rangle \to (+\hbar,0,-\hbar)\ \text{的概率分别为}\ \frac14,\frac12,\frac14;
\]
\[
|0\rangle \to (+\hbar,0,-\hbar)\ \text{的概率分别为}\ \frac12,0,\frac12;
\]
\[
|-1\rangle \to (+\hbar,0,-\hbar)\ \text{的概率分别为}\ \frac14,\frac12,\frac14.
\]
故全部可能的结果组合及其概率为
\ansmath{\begin{aligned}
&(S_z=+\hbar,\ S_y=+\hbar), &&P=\frac14\cos^4\frac{Dt}{2},\\
&(S_z=+\hbar,\ S_y=0), &&P=\frac12\cos^4\frac{Dt}{2},\\
&(S_z=+\hbar,\ S_y=-\hbar), &&P=\frac14\cos^4\frac{Dt}{2},\\
&(S_z=0,\ S_y=+\hbar), &&P=\frac14\sin^2 Dt,\\
&(S_z=0,\ S_y=-\hbar), &&P=\frac14\sin^2 Dt,\\
&(S_z=-\hbar,\ S_y=+\hbar), &&P=\frac14\sin^4\frac{Dt}{2},\\
&(S_z=-\hbar,\ S_y=0), &&P=\frac12\sin^4\frac{Dt}{2},\\
&(S_z=-\hbar,\ S_y=-\hbar), &&P=\frac14\sin^4\frac{Dt}{2}.
\end{aligned}}
其中 $(S_z=0,S_y=0)$ 的概率为 $0$。
\end{solution}
\end{question}

\begin{question}
考虑圆环上的两个无相互作用、无自旋全同粒子。单粒子哈密顿量为
\[
\hat h=-\frac{\hbar^2}{2mR^2}\frac{d^2}{d\phi^2},
\]
单粒子归一化本征态为
\[
\varphi_n(\phi)=\frac{1}{\sqrt{2\pi}}e^{in\phi},
\qquad
\varepsilon_n=\frac{\hbar^2 n^2}{2mR^2},
\qquad n\in\mathbb Z.
\]
总哈密顿量为
\[
\hat H=\hat h_1+\hat h_2.
\]

\begin{enumerate}[label=(\arabic*),leftmargin=2.5em]
  \item 对两个全同玻色子，写出全部能级与归一化本征态；
  \item 对两个全同费米子，写出全部能级与归一化本征态；
  \item 求玻色子基态的期望值 $\langle \cos(\phi_1-\phi_2)\rangle$；
  \item 取一个费米子基态
  \[
  \Psi_{0,1}^{(F)}(\phi_1,\phi_2)=\frac{\varphi_0(\phi_1)\varphi_1(\phi_2)-\varphi_1(\phi_1)\varphi_0(\phi_2)}{\sqrt2},
  \]
  求该态中的 $\langle \cos(\phi_1-\phi_2)\rangle$。
\end{enumerate}

\begin{solution}
两个玻色子的总波函数必须对称，因此当 $n_1=n_2$ 时
\[
\Psi_{n,n}^{(B)}(\phi_1,\phi_2)=\varphi_n(\phi_1)\varphi_n(\phi_2),
\qquad
E_{n,n}=2\varepsilon_n.
\]
当 $n_1<n_2$ 时
\ansmath{\Psi_{n_1,n_2}^{(B)}=\frac{\varphi_{n_1}(\phi_1)\varphi_{n_2}(\phi_2)+\varphi_{n_2}(\phi_1)\varphi_{n_1}(\phi_2)}{\sqrt2},
\qquad
E_{n_1,n_2}=\varepsilon_{n_1}+\varepsilon_{n_2}.}
这给出了全部玻色子能级和归一化本征态。

两个费米子的总波函数必须反对称，因此只能有 $n_1<n_2$：
\ansmath{\Psi_{n_1,n_2}^{(F)}=\frac{\varphi_{n_1}(\phi_1)\varphi_{n_2}(\phi_2)-\varphi_{n_2}(\phi_1)\varphi_{n_1}(\phi_2)}{\sqrt2},
\qquad
E_{n_1,n_2}=\varepsilon_{n_1}+\varepsilon_{n_2}.}

玻色子基态为两粒子都占据 $n=0$：
\[
\Psi_{0,0}^{(B)}=\frac{1}{2\pi}.
\]
于是
\[
\langle \cos(\phi_1-\phi_2)\rangle
=\frac{1}{4\pi^2}\int_0^{2\pi}d\phi_1\int_0^{2\pi}d\phi_2\,\cos(\phi_1-\phi_2)=0.
\]
故
\ansmath{\langle \cos(\phi_1-\phi_2)\rangle_{\Psi_{0,0}^{(B)}}=0.}

对费米子态 $\Psi_{0,1}^{(F)}$，其概率密度为
\[
|\Psi_{0,1}^{(F)}|^2=\frac{1}{4\pi^2}\bigl(1-\cos(\phi_1-\phi_2)\bigr).
\]
因此
\[
\langle \cos(\phi_1-\phi_2)\rangle
=\frac{1}{4\pi^2}\int d\phi_1d\phi_2\,\cos(\phi_1-\phi_2)\bigl(1-\cos(\phi_1-\phi_2)\bigr).
\]
第一项积分为零，第二项给出 $-1/2$，故
\ansmath{\langle \cos(\phi_1-\phi_2)\rangle_{\Psi_{0,1}^{(F)}}=-\frac12.}
\end{solution}
\end{question}

\begin{question}
考虑一维对称有限深势阱
\[
V(x)=
\begin{cases}
-V_0,& |x|<a,\\
0,& |x|>a,
\end{cases}
\qquad V_0>0.
\]
考虑束缚态 $-V_0<E<0$。

\begin{enumerate}[label=(\arabic*),leftmargin=2.5em]
  \item 写出偶宇称与奇宇称束缚态的波函数形式，并导出它们各自满足的本征值方程；
  \item 证明该势阱总至少存在一个偶宇称束缚态；
  \item 证明存在奇宇称束缚态的条件为
  \[
  a\sqrt{\frac{2mV_0}{\hbar^2}}>\frac{\pi}{2};
  \]
  \item 在深阱极限下，写出前几个束缚能级的近似表达式。
\end{enumerate}

\begin{solution}
记
\[
k=\frac{\sqrt{2m(E+V_0)}}{\hbar},
\qquad
\kappa=\frac{\sqrt{-2mE}}{\hbar},
\qquad
k^2+\kappa^2=\frac{2mV_0}{\hbar^2}.
\]

偶宇称束缚态可写为
\[
\psi(x)=
\begin{cases}
A\cos(kx),& |x|<a,\\
B e^{-\kappa |x|},& |x|>a.
\end{cases}
\]
在 $x=a$ 处连续并令导数连续，得
\ansmath{k\tan(ka)=\kappa.}

奇宇称束缚态可写为
\[
\psi(x)=
\begin{cases}
A\sin(kx),& |x|<a,\\
B\,\mathrm{sgn}(x)e^{-\kappa |x|},& |x|>a.
\end{cases}
\]
同样由边界匹配得到
\ansmath{-k\cot(ka)=\kappa.}

对偶宇称方程，函数 $y=k\tan(ka)$ 在区间 $0<ka<\pi/2$ 上从 $0$ 连续增至 $+\infty$；而 $\kappa=\sqrt{2mV_0/\hbar^2-k^2}$ 在该区间上是正的连续函数。因此两条曲线必有交点，所以总至少存在一个偶宇称束缚态。

对奇宇称方程，要有解必须使得 $ka$ 能进入区间 $(\pi/2,\pi)$。由于 $ka$ 的最大可能值为
\[
ka\le a\sqrt{\frac{2mV_0}{\hbar^2}},
\]
所以存在奇宇称束缚态的必要且充分条件是
\ansmath{a\sqrt{\frac{2mV_0}{\hbar^2}}>\frac{\pi}{2}.}

在深阱极限下，$\kappa a\gg 1$，阱壁近似为无限高，故宽度为 $2a$ 的无限深势阱近似成立。于是能级近似为
\ansmath{E_n\approx -V_0+\frac{\hbar^2\pi^2 n^2}{8ma^2},
\qquad n=1,2,3,\dots.}
其中 $n=1,3,5,\dots$ 对应偶宇称，$n=2,4,6,\dots$ 对应奇宇称。
\end{solution}
\end{question}

